\section{The v0.5 mechanisms}
\label{sec:mechanisms}

\vfive{} is a new policy registered alongside \vfour{}, which stays unchanged and
remains the reference opponent. Every mechanism is individually switchable, so
each can be ablated; the switch names are given with each mechanism. Three are
built and measured (\S\ref{sec:mech-built}); seven are specified and not built
(\S\ref{sec:mech-designed}); four plausible repairs were measured, rejected, and
are recorded here so that they are not rebuilt (\S\ref{sec:mech-rejected}).

\subsection{Built: M1, M2 and M8}
\label{sec:mech-built}

\paragraph{M1 --- live-ask gating and ownership scaling (\texttt{m1}, \texttt{m1p}).}
The candidate set is restricted to asks that are not provably dead in the sense of
Definition~\ref{def:dead}, computed from the actor's own deduction state. Because
the predicate is a hard deduction rather than a posterior quantity, the filter
cannot introduce a modelling error; it removes moves whose probability of success
is exactly zero and nothing else. When the filtered set is empty --- a genuinely
starved turn, \numStarved\% of asks --- the policy falls back to the full legal
set, because the rules still oblige it to move. A second half was designed
alongside it and is \emph{not} in the shipped configuration: under \texttt{m1p},
all four ownership features of \S\ref{sec:diag-features} --- own-set progress,
team control, lock completion and the completion bonus --- are multiplied by the
hit probability, so that a card nobody can hand over cannot outscore a live one
inside the surviving set. The two halves are separately switchable because they
are separate claims: the gate changes which asks are available, the scaling
changes how the survivors are ranked. The gate ships on
(\texttt{liveAskGate}); the scaling ships \emph{off}
(\texttt{ownershipByP}, false in \filepath{engine/src/v05.hpp}), because it was
measured and rejected --- once the gate has removed the $p = 0$ case there is no
dead ask left for the scaling to suppress, and what remains costs
\numMonePReject{} points on top of the other mechanisms at design time and
$\numAblOwnershipPDelta$ at the frozen configuration
(\S\ref{sec:results-ablations}). Every claim in this report about \vfive{}'s play
is a claim about the gate alone.

\paragraph{M2 --- capacity-feasible allocation (\texttt{m2}).} The per-card argmax
of \S\ref{sec:diag-forced} is replaced by a search over allocations that respect
the declarer's own capacity constraints. A half-suit has six cards and a team has
three seats, so the assignment space is at most $3^6 = 729$; the implementation
enumerates it, rejects any assignment that places a card at a seat excluded by a
certificate or that gives a teammate more cards than their remaining hand can
hold, and scores the survivors with the same joint estimator the declaration rule
already uses --- which runs on the deployed approximate belief, so the resulting
probability is an estimate and not an exact posterior value. The search is
exhaustive over the feasible set and is cheaper than the dynamic program it
replaces. The same routine feeds three call sites: the forced-endgame
declaration, the willingness ladder that selects the forced declarer --- which
recovers its function for free, since the statistic it thresholds is no longer
identically zero --- and, under a separate switch, the voluntary declaration
path.

\paragraph{M8 --- termination without misdeclaration (\texttt{stage2},
\texttt{norepeat}).} The second stage of the event-count escalation is off by
default; the first stage, which requires a better-than-even allocation
probability, is the only escalation that remains. With M1 in place the deadlock
that the second stage existed to break is gone. The design intended to put a
structural backstop where the event-count guillotine had been --- a repetition
guard preventing the same agent from asking the same (card, target) pair twice,
applied inside the candidate filter and relaxed before the starved-turn fallback
--- and that guard is implemented, switchable, and \emph{off} in the shipped
configuration (\texttt{repeatGuard}, false in \filepath{engine/src/v05.hpp}): it
costs \numGuardReject{} points against \vfour{} at design time and
$\numAblRepeatGuardDelta$ at the frozen configuration, for the reason
\S\ref{sec:results-ablations} gives. M8 as shipped is therefore a deletion and
nothing else. The design rule behind it is the project owner's, and it is a claim
about the game: a risky ask is better than a declaration the policy itself scores
at zero.

\paragraph{What each is worth alone.} Each mechanism was measured in isolation
against \vfour{} over \numMechDeals{} deals in six seat rotations at seed
\numMechSeed{}, with \vfive{} configured to reproduce \vfour{} exactly as the
control (\numControlWin\%). The results are in Table~\ref{tab:mechanisms}. Two
readings are worth stating plainly. First, the three mechanisms that remove the
pathology are close to strength-neutral in head-to-head play, and the one that is
not --- M8 --- gains by \emph{withholding} a bad declaration rather than by
playing better. Second, the repetition guard on its own is exactly neutral, which
is consistent with \S\ref{sec:diag-features}: with the dead asks already gone,
there is almost nothing left for it to catch.

\begin{table}[t]
\centering
\caption{Single-mechanism win rate against \vfour{}, \numMechDeals{} deals in six
rotations, seed \numMechSeed{}. The control is \vfive{} with every mechanism
disabled, which reproduces \vfour{}. These are single-opponent measurements and
are not a strength claim for the assembled policy; see \S\ref{sec:results}.}
\label{tab:mechanisms}
\begin{tabular}{llr}
\toprule
\textbf{Configuration} & \textbf{Mechanism} & \textbf{win rate} \\
\midrule
control (all off)       & ---                                & \numControlWin\% \\
\texttt{m1}             & live-ask gate only                 & \numMoneAlone\% \\
\texttt{m1p}            & ownership features scaled by $p$   & \numMonePAlone\% \\
\texttt{m2}             & capacity-feasible allocation       & \numMtwoAlone\% \\
\texttt{stage2=0}       & escalation stage 2 removed         & \numMeightAlone\% \\
\texttt{norepeat}       & repetition guard only              & \numGuardAlone\% \\
\bottomrule
\end{tabular}
\end{table}

\paragraph{The pathology gates.} The build order was gated on the three pathology
statistics rather than on win rate, following the standing preference that
robustness is the research question. \S\ref{sec:results} reports the full
comparison; the gates themselves were: after M1, no game may contain a dead run
of six; after M2, forced-endgame accuracy must leave zero; after M8, no game may
reach the action cap and no declaration may be made past the horizon.

\paragraph{The parameters behind this table are not \vfive{}'s.}
Table~\ref{tab:mechanisms} was measured before the refit, with the
\numFitDims{} fitted coordinates still \vfour{}'s, and on the design-time
assembly, which carried both of the mechanisms this subsection has just recorded
as rejected. They had been fitted with the escalation of
\S\ref{sec:diag-dilemma} present as a safety net, so with that net removed that
assembly under-claims: it makes \numVfourParamDecl{}
declarations per game against \vfour{}'s \numVfourDecl{}, and scores
\numVfourParamWin\% [\numVfourParamCI] head-to-head. Most of that shortfall
belongs to the two rejected mechanisms rather than to the coordinates, and the
distinction matters because the figure has been read as a parameter mismatch:
on the same bank, with the same \vfour{} coordinates, the assembly that actually
ships scores \numShipParamSameBank\% [\numShipParamSameCI], which is
\numRejectedCost{} points higher and level with \vfour{}
(\filepath{research/v05/runs/v04vector-in-v05.txt}).
\S\ref{sec:fitting} describes the refit that replaced those
coordinates in the frozen configuration, and every strength figure for \vfive{}
in this report is measured after it; the ablations of
\S\ref{sec:results-ablations} repeat this table's comparisons on the fitted
vector, and its control is correspondingly not \vfour{}.

\subsection{Designed and not built}
\label{sec:mech-designed}

None of the following is on \vfive{}'s decision path. Prototype code exists
outside it for the target-scoring and opponent-model mechanisms, and is used only
to check that the quantities they need can be computed; nothing below has been
measured in play, and each entry states the constraint the diagnosis places on
its eventual construction.

\paragraph{M3 --- a net-information term in the ask score.} Replace the negative
value-of-information term with a signed net term: the information the ask
delivers to teammates, minus the information it delivers to opponents, weighted
separately. This is the secrecy-rate trade-off of \S\ref{sec:related-info}, and
the two filters the engine already produces --- the policy-agnostic posterior and
the policy-weighted one --- are enough to compute both sides. The design
constraint established in \S\ref{sec:diag-information} is that the channel is
almost never \emph{provably} free: it must be priced against the probability that
the team owns the half-suit, which at the relevant states averages
\numOwnProbMean{} and peaks at \numOwnProbMax{}, not
gated on a certainty no observer has. A \vfive{} that signals only when signalling
is free will never signal. M3 requires M4.

\paragraph{M4 --- a knowledge model of the other five seats.} Maintain, for each
seat, the public deduction state plus a posterior over that seat's hand. Nearly
free because the transcript is public, and a precondition for M3, M5 and M7.

\paragraph{M5 --- target-dimension selection.} Score the target explicitly:
lockout value (which opponent must not receive the turn), void progress inside
the half-suit, and, where conventions are enabled, codebook value. The measured
headroom is in \S\ref{sec:diag-channels}, and the verification recorded there
requires that the choice be priced against the capacity-marginal hit probability
rather than treated as free.

\paragraph{M6 --- partner-aware stochastic action selection.}
\label{sec:mech-partner}
Select by softmax over the near-optimal set rather than by argmax. This is
required for any leak penalty to mean anything at all: leakage is a function of
the distribution over actions, so a deterministic policy sits at the
deterministic bound whatever its feature weights \cite{alvim-leakage}. The
temperature and the convention flag are both keyed on the partner regime. With
bot teammates --- the self-play condition, in which partners follow a known
blueprint --- lower temperature and convention-rich play are appropriate, because
the receiver decodes what the sender encodes. With human or unknown teammates the
blueprint assumption fails, and play must be grounded and less predictable,
because an opponent who has watched a deterministic argmax across many games
learns to read it. The two regimes are to be reported separately, and the
self-play configuration is never to be headlined as though it were the one a
human will meet.

\paragraph{M7 --- an online per-seat opponent model.} Replace the two global
scalars with a per-player posterior over a small type library, updated from that
player's public asks and from a \emph{time-local} silence statistic --- own turns
since last asking in this half-suit, and having been asked in it without
replying --- rather than from a whole-game ask count, which
\S\ref{sec:diag-prior} shows is inert. Shape it as a data-biased response
\cite{johanson-dbr}: confidence per information set, decaying to the population
prior where evidence is absent, so that a deliberately silent opponent produces
no evidence and receives the prior instead of being misread. Expose the
exploitation/robustness dial explicitly and report the frontier
\cite{johanson-rnash} rather than a single operating point.

\paragraph{M9 --- the value model, rebuilt or retired.} Two options, to be decided
by measurement: collect rows at declaration points as well as ask points, drop
the degenerate features, and refit jointly with the policy vector so that the
freeze step covers it; or retire it and score declarations directly on the
allocation probability against a fitted threshold. \S\ref{sec:diag-value} shows
that the ask-side use contributes nothing and the declaration-side use does, and
that a tree model finds real structure the linear form cannot express.

\paragraph{M10 --- fitting and evaluation.} Put the mirror in the panel. Replace
the soft minimum, which at the panel's spread is a weighted mean
(\S\ref{sec:diag-value}), with a worst-case or minimax-regret objective over the
style set, and report both the minimum and the regret --- a raw minimum over a
small style set stagnates on the hardest member \cite{beukman-remidi}. Promote
the three pathology statistics --- longest dead run, share of games with a dead
run of six, share of post-horizon declarations that are wrong --- to first-class
metrics gating every commit.

\subsection{Conventions, and which configuration is the headline}
\label{sec:mech-conventions}

The strategy corpus states that pre-agreed conventions are outside this game
\cite{develin}; the rules corpus documents a named convention in real use
\cite{pagat}; and a pair of policies trained together has an implicit
pre-agreement whether or not anyone writes it down. \vfive{} therefore builds the
signalling machinery behind an explicit switch, ships it \emph{off} as the
headline configuration --- which is legal under every reading of the rules we
found --- and publishes the with/without difference as a result in its own right
rather than folding it into a single number.

\subsection{Measured, rejected, and recorded so as not to be rebuilt}
\label{sec:mech-rejected}

\begin{enumerate}
  \item \textbf{A time-varying holding cost in the value model.} The literature
        prescription of \S\ref{sec:related-stopping} --- make waiting strictly
        costly and the stopping problem becomes well posed
        \cite{mguni-actioncost} --- was implemented as a linear penalty in
        elapsed events and swept over a forty-fold range of magnitudes. It is a
        mean-shifter. Events per game fall from \numHoldEventsBefore{} to
        \numHoldEventsAfter{} and the frequency of dead runs from
        \numHoldRunGamesBefore\% to \numHoldRunGamesAfter\%, but the tail does
        not move at all: the 90th percentile of events goes from
        \numHoldPninetyBefore{} to \numHoldPninetyAfter{} and the longest dead
        run from \numHoldLongestBefore{} to \numHoldLongestAfter{}. Declaration
        error rises from \numHoldDeclBefore\% to \numHoldDeclAfter\% and the
        configuration loses \numHoldLoss{} points head-to-head. A constant holding cost is
        also not missing --- the declaration margin already is one --- and in
        deadlocked positions the urgency flag bypasses the value-based
        declaration rule entirely, so the term binds in fewer than one in ten
        late blocked opportunities.
  \item \textbf{Deleting the policy prior.} Worse, not better: by
        \numAgnosticCost{} points [\numAgnosticCI] on a mixed panel of three
        deceptive archetypes and the honest mirror, and by
        \numPriorDeleteCost{} points [\numPriorDeleteCI] against the deceptive
        archetypes alone (\S\ref{sec:diag-prior}). The exposure is
        over-weighting.
  \item \textbf{A turn-transfer willingness ladder.} A genuine multi-candidate
        transfer decision arises \numTransferRate{} times per game and governs
        \numTransferAskShare\% of asks; the unilateral choice already matches an omniscient
        oracle at \numTransferOracle{} of \numTransferOracleN{} of them, and
        handing one team a ground-truth chooser is worth
        $\numTransferValue \pm \numTransferSE$ sets per game over \numTransferGames{}
        paired games --- indistinguishable from zero. The channel is also the leakier of
        the two: an eight-rung ladder transmits \numLadderLeak{} bits per
        candidate against the one bit the rules license, and it transmits them
        mid-game to live opponents. Budget nothing for it.
  \item \textbf{Confidence-ranked declaration arbitration.} Resolving
        simultaneous declarations by lowest seat rather than by confidence costs
        \numArbCost{} percentage points pooled over \numArbGames{} games. The
        rules-legal willingness ladder may be added if it is convenient, but
        nothing should be engineered for it.
\end{enumerate}
